%% kapa: exam=1 round=23 year=2012 period=2 form=A subject=appraisal_law subject_ko=부동산관계법규 seq=37 number=77
%% source: Q-Net 감정평가사 기출문제 (cst00309) artlSeq=5052580 file=raw/1008/5052580/2012년 제23회 감정평가사 제1차시험 기출문제(2교시 A형).pdf
%% conversion: kordoc 4.24.51 -> tools/gen1.py (auto); review=proofread-sonnet; flags=none
\begin{question}{77}
부동산등기법상 등기에 관한 설명으로 옳지 않은 것은?
\begin{choices}
  \item 등기는 법률에 다른 규정이 없는 경우에는 등기권리자와 등기의무자가 공동으로 신청한다.
  \item 부동산표시의 변경등기는 소유권의 등기명의인이 단독으로 신청한다.
  \item 판결에 의한 등기는 승소한 등기권리자 또는 등기의무자가 단독으로 신청한다.
  \item 소유권보존등기의 말소등기는 등기명의인이 단독으로 신청한다.
  \item 등기관이 등기를 마친 경우 그 때부터 등기는 효력을 발생한다.
\end{choices}
\end{question}
\answer{5}
